\documentclass[12pt, a4]{article}

\usepackage{amsmath}
\usepackage{german}
\usepackage{graphicx}
\DeclareMathOperator\sgn{sgn}

%\graphicspath{{/home/ritz/Teaching/fcn/fig/}}

%----------------- layout ----------------------------------------

\oddsidemargin=0in
\textwidth=6in
\topmargin=-0.7in
\textheight=9.5in
\baselineskip=16pt
\font\mittel=cmr17
\font\roman=cmr12
\font\klein=cmr9
\font\mgross=cmr10
\newfont{\CAP}{cmcsc10 scaled 900}
\pagestyle{empty}
\parindent=0mm

%----------------- macros ----------------------------------------
\newcommand{\conc}[1]{[{\rm C}]_{\text{#1}}}
\newcommand{\units}[1]{\text{#1}}
\newcommand{\kalium}{$\text{K}^+$}
\newcommand{\natrium}{$\text{Na}^+$}
\newcommand{\chlor}{$\text{Cl}^-$}
\newcommand{\wasser}{$\text{H}_2\text{O}$}
%-----------------------------------------------------------------


\begin{document}

\parbox{2cm}{
\includegraphics[width=1.8cm]{bccnlogo.pdf}
}
\parbox{11cm}{
\begin{center}
\large HUMBOLDT-UNIVERSIT"AT \hskip 0.1 cm ZU \hskip 0.1 cm BERLIN
\vskip 0.1 true cm
\mgross BERNSTEIN CENTER FOR COMPUTATIONAL NEUROSCIENCE
\end{center}
}
\parbox{2cm}
{
\hfill
\includegraphics[width=1.8cm]{hublogo.pdf}
}

\vskip 0.6 true cm

\leftline{\CAP Humboldt-Universit"at zu Berlin 
\hfill Phone: 030/2093-9110}
\leftline{\CAP Philippstr. 13 House 6
              \hfill Fax: 030/2093-6771}
\leftline{\CAP \hfill webpage: http://www.bccn-berlin.de/}


\vskip 0.8 true cm
\centerline{\bf Models of Neural Systems, WS 2010/11}
\centerline{\bf Computer Practical 4}
\centerline{Solutions to hand in on: November, 22nd, 2010}

\vskip 0.8 true cm
{\bf Visual receptive fields -- exercises:}
\begin{enumerate}

    \item \textbf{Visual stimulus.} A commonly used stimulus for the recordings from the visual
        system is a sinusoidal grating:
        \begin{equation}
            s(x, y) = A \cos(Kx\cos\Theta + Ky\sin\Theta
            -\Phi),
            \label{eq:grating}
        \end{equation}
        where $x$, $y$ are the spatial coordinates, $K$ is the spatial
        frequency, $\Theta$ is the orientation,  $\Phi$ is the spatial
        phase and $A$ is the contrast amplitude.

        \begin{enumerate}
            \item Approximate the visual field $(x,y)$ with a
                2-dimensional grid of uniformly-distributed
                photoreceptors (retina):
                $(x_i, y_i)= (-x_0+i\Delta x, -y_0+j\Delta y)$ for
                $i=0,1,\dots, \frac{2x_0}{\Delta x}$ and $j=0,1,\dots,
                \frac{2y_0}{\Delta y}$ where $\Delta x$, $\Delta y$
                are bin sizes and $x_0$ and $y_0$ are the
                visual field extents (both in degrees). Take
                $x_0=5^\circ$, $y_0=5^\circ$.
                
                \textit{Hint:} You can use
                \texttt{pylab.meshgrid} to generate such a grid.
                % explain meshgrid -- people were confused!
            \item \label{ex:grating} Plot the grating on the
                grid for the following parameters: $\Theta=0$,
                $\Phi=0$, $K=\pi\,\frac{1}{\mathrm{degree}}$, $A=1$.  Use the 
                \texttt{pylab.imshow} function to show a 2-dimensional map. 
            \item \label{ex:varygrating}Vary the orientation $\Theta$ and the
                spatial frequency $K$ and plot the resulting maps.
        \end{enumerate}

    \item \textbf{A model of receptive fields}.  A receptive field (RF)
                describes the ``sensitivity'' of a neuron to inputs arriving from
                different points of the visual field. A convenient mathematical
                approximation of a simple cell receptive field is provided by a
                Gabor function:

        \begin{equation}
           D_s(x,y)=\frac{1}{2\pi\sigma_x\sigma_y}
            \exp\left(-\frac{x^2}{2\sigma_x^2}-\frac
            {y^2}{ 2\sigma_y^2}\right)\cos(kx-\phi),
            \label{eq:Gabor}
        \end{equation}
        where  $k$ is the preferred spatial frequency, $\phi$ is
        the preferred spatial phase,  $\sigma_x$,
        $\sigma_y$ are the receptive field sizes.

        \begin{enumerate}
            \item \label{ex:Gabor}  Implement a Gabor function as a model of
                a receptive field with preferred spatial frequency
                $k=\pi\,\frac{1}{\mathrm{degree}}$,
                preferred spatial phase $\phi=0$
                and receptive field size $\sigma_x=2^\circ$,
                $\sigma_y=1^\circ$. Plot this Gabor receptive field.
                
            \item \label{ex:RF_questions} What are the optimal
                stimuli for a neuron with such a RF? Where would you
                expect to find a neuron with such a RF?  Why does this
                RF characterize a simple cell but not a complex cell?
                
            \item \label{ex:linear_response} In order to calculate the
                response of a neuron to a visual stimulus, we need to
                sum the input intensity to all photoreceptors in the
                visual field weighted by the receptive field.
                Mathematically this can be represented as a linear
                filter:

                \begin{equation}
                    L_s = \sum_{i=0}^{\frac{2x_0}{\Delta x}}
                    \sum_{j=0}^{\frac{2y_0}{\Delta y}} D_s(x_i,y_i) s(x_i,y_i)
                    \Delta x \Delta y.
                    \label{eq:respons}
                \end{equation}
                
                Numerically calculate the response of a cell with a
                Gabor receptive field (Equation \ref{eq:Gabor},
                parameters as in exercise \ref{ex:Gabor}) to the
                grating from exercise \ref{ex:grating}, and state the
                resulting number in your report. Compare it with a
                response to an uniformly illuminated visual field
                ($s(x,y)=A$). How would you explain the difference?

        \end{enumerate}

    \item\label{ex:tune_curve} \textbf{Tuning Curve}. A tuning curve describes the dependence of
        the neuronal response on a parameter of the stimulus such as
        spatial frequency, orientation or phase. Plot the neuronal
        response $L_s$ to the grating from the exercise \ref{ex:grating} as a 
        function of:

        \begin{enumerate}
            \item\label{ex:tuneorient} stimulus orientation $\Theta$,
            \item ratio of the stimulus' spatial frequency to the
                cell's preferred value $K/k$,
            \item\label{ex:tunephase} stimulus' spatial phase  $\Phi$.

        \end{enumerate}
        Characterize the receptive field in terms of its selectivity
        (spatial, orientation and spatial frequency) and preferred stimulus
        (the stimulus giving rise to maximal response). 
    \item {\bf Report.} Submit plain source code files for all exercises and a
    report. The report should be in PDF format and include source
    code to all exercises, answers to the
    questions in exercises \ref{ex:RF_questions}, \ref{ex:linear_response}
    and \ref{ex:tune_curve}, and plots from exercises~\ref{ex:grating},
    \ref{ex:varygrating}, \ref{ex:Gabor} and~\ref{ex:tune_curve} with
    caption and labels. 

\end{enumerate}

{\bf Reference:}

D. H. Hubel and T. N. Wiesel (1959). Receptive fields of single
neurones in the cat's striate cortex, J Physiol. 148(3): 574-591. 

\vfill
\centerline{\CAP Contact}
\CAP

\begin{tabular}{lll}
Urs Bergmann & Phone: 2093-8924 & Email:
urs.bergmann(at)biologie.hu-berlin.de \\
Matthias Guggenmos & & Email: matthias.guggenmos(at)bccn-berlin.de \\
Richard Kempter \hfill & Phone: 2093-8925 \hfill & Email:
r.kempter(at)biologie.hu-berlin.de \\
Robert Schmidt & Phone: 2093-8926 & Email: r.schmidt(at)biologie.hu-berlin.de
\\
Bartosz Telenczuk & Phone: 2093-8838 & Email:
b.telenczuk(at)biologie.hu-berlin.de \\
\end{tabular}

\end{document}
 
